\documentclass[11pt]{article}

\usepackage[margin=1in]{geometry}
\usepackage[T1]{fontenc}
\usepackage{lmodern}

\usepackage{amsmath, amssymb}

\usepackage{setspace}

\usepackage[most]{tcolorbox}
\usepackage{xcolor}

\tcbset{
  enhanced,
  boxrule=0pt,
  frame hidden,
  breakable
}

\usepackage{fancyhdr}
\pagestyle{fancy}
\fancyhf{}
\setlength{\headheight}{52pt}
\renewcommand{\headrulewidth}{0pt}

\newcommand{\Course}{Math 4610 Spring 2026}
\newcommand{\Instructor}{Homework 2}
\newcommand{\AssignmentType}{Homework} % or Discussion
\newcommand{\AssignmentNumber}{2}

\newcommand{\Contributors}{
Marks, Stephen
}

\fancyhead[L]{\Course\\ \Instructor}
\fancyhead[R]{\Contributors}

\newcommand{\ProblemDivider}[1]{%
  \vspace{1.0em}
  \noindent\textbf{#1}\par
  \vspace{0.35em}
}

\definecolor{problemBack}{RGB}{235,242,250}
\definecolor{problemFrame}{RGB}{70,110,160}
\definecolor{exerciseGray}{gray}{0.96}
\definecolor{exerciseGrayFrame}{gray}{0.45}

\newtcolorbox{problemBox}{
  colback=exerciseGray,
  colframe=exerciseGrayFrame,
  arc=2mm,
  left=8pt,right=8pt,top=8pt,bottom=8pt
}

\newtcolorbox{solutionBox}{
  colback=white,
  colframe=white,
  arc=0mm,
  left=0pt,right=0pt,top=0pt,bottom=0pt
}

\newcommand{\ProblemAnnounce}{\noindent}
\newcommand{\SolutionAnnounce}{\noindent\textbf{Solution.}\par\medskip}

\newcommand{\Exercise}[1]{\textbf{Exercise~#1.}\;}

\newcommand{\R}{\mathbb{R}}
\newcommand{\C}{\mathbb{C}}
\newcommand{\F}{\mathbb{F}}
\newcommand{\Lmap}{\mathcal{L}}
\newcommand{\inv}{^{-1}}

\begin{document}

\begin{center}
  {\Large \textbf{\AssignmentType~\AssignmentNumber}}
\end{center}
\vspace{0.75em}

\ProblemDivider{1}

\begin{problemBox}
\ProblemAnnounce
\Exercise{5.A.3}
Suppose \( T \in \Lmap(V) \). Prove that the intersection of every collection of
subspaces of \( V \) invariant under \( T \) is invariant under \( T \).
\end{problemBox}

\begin{solutionBox}
\SolutionAnnounce
\doublespacing

Let \(\mathcal{A} = \{U_i\}_{i \in I}\) be a collection of subspaces of \(V\) each invariant under \(T\). 
Let \(W = \bigcap_{i \in I} U_i\). Since each \(U_i\) is a subspace, \(W\) is a subspace of \(V\) (by 1.C.11). 

Take any \(w \in W\). Then \(w \in U_i\) for every \(i \in I\). Because each \(U_i\) is invariant under \(T\), 
\(T(w) \in U_i\) for every \(i\). So \(T(w) \in \bigcap_{i \in I} U_i = W\).

Thus \(W\) is invariant under \(T\).

\end{solutionBox}

\ProblemDivider{2}

\begin{problemBox}
\ProblemAnnounce
\Exercise{5.A.5}
Suppose \( T \in \Lmap(\R^2) \) is defined by \( T(x, y) = (-3y, x) \). Find the eigenvalues of \( T \).
\end{problemBox}

\begin{solutionBox}
\SolutionAnnounce
\doublespacing

Suppose \(\lambda \in \R\) is an eigenvalue of \(T\) with eigenvector \((x,y) \neq (0,0)\). Then
\[
T(x,y) = \lambda (x,y) \quad \Rightarrow \quad (-3y, x) = (\lambda x, \lambda y).
\]
This yields the system
\[
\begin{cases}
-3y = \lambda x,\\
x = \lambda y.
\end{cases}
\]
Substituting \(x = \lambda y\) into the first equation gives
\[
-3y = \lambda (\lambda y) = \lambda^2 y.
\]
If \(y = 0\), then \(x = 0\) from the second equation, contradicting \((x,y) \neq (0,0)\) (by 5.8). Thus \(y \neq 0\), and we may divide by \(y\) to obtain
\[
\lambda^2 = -3.
\]
No real \(\lambda\) satisfies this equation. Therefore \(T\) has no (real) eigenvalues.

\end{solutionBox}

\ProblemDivider{3}

\begin{problemBox}
\ProblemAnnounce
\Exercise{5.A.11}
Suppose \(V\) is finite-dimensional, \(T \in \Lmap(V)\), and \(\alpha \in \F\). Prove that
there exists \(\delta > 0\) such that \(T - \lambda I\) is invertible for all \(\lambda \in \F\) with \(0 < |\alpha - \lambda| < \delta\).
\end{problemBox}

\begin{solutionBox}
\SolutionAnnounce
\doublespacing

If \(T\) has no eigenvalues, then by Theorem 5.7 (or the analogous result in the text) \(T - \lambda I\) is invertible for every \(\lambda \in \F\); we may choose any \(\delta > 0\).

Now suppose \(T\) has eigenvalues. Because \(V\) is finite‑dimensional, \(T\) has only finitely many distinct eigenvalues; list them as \(\lambda_1, \dots, \lambda_n\).

If \(\alpha\) equals one of the \(\lambda_k\), let
\[
\delta = \min\{|\alpha - \lambda_j| : \lambda_j \neq \alpha\}.
\]
If \(\alpha\) is not an eigenvalue, let
\[
\delta = \min\{|\alpha - \lambda_j| : j = 1,\dots,n\}.
\]
In either case, \(\delta > 0\) (the minimum is taken over a finite set of positive numbers).

For any \(\lambda \in \F\) with \(0 < |\alpha - \lambda| < \delta\), the distance from \(\lambda\) to every eigenvalue \(\lambda_j\) is positive, hence \(\lambda\) is not an eigenvalue of \(T\). By Theorem 5.7, \(T - \lambda I\) is invertible.

\end{solutionBox}

\ProblemDivider{4}

\begin{problemBox}
\ProblemAnnounce
\Exercise{5.A.13}
Suppose \(T \in \Lmap(V)\) and \(S \in \Lmap(V)\) is invertible.
\begin{enumerate}
  \item[(a)] Prove that \(T\) and \(S^{-1}TS\) have the same eigenvalues.
  \item[(b)] What is the relationship between the eigenvectors of \(T\) and the eigenvectors of \(S^{-1}TS\)?
\end{enumerate}
\end{problemBox}

\begin{solutionBox}
\SolutionAnnounce
\doublespacing

\begin{enumerate}
  \item[(a)] Let \(\lambda \in \F\).  
  \(\lambda\) is an eigenvalue of \(T\)  
  \(\iff\) there exists \(u \in V\), \(u \neq 0\), such that \(Tu = \lambda u\)  
  \(\iff\) there exists \(u \in V\), \(u \neq 0\), such that \(S^{-1}TS (S^{-1}u) = \lambda (S^{-1}u)\) (set \(v = S^{-1}u\), which is nonzero because \(S\) is invertible)  
  \(\iff\) \(\lambda\) is an eigenvalue of \(S^{-1}TS\).

  \item[(b)] The correspondence is given by \(v \mapsto Sv\) from eigenvectors of \(S^{-1}TS\) to eigenvectors of \(T\), and \(u \mapsto S^{-1}u\) from eigenvectors of \(T\) to eigenvectors of \(S^{-1}TS\).  
  So,
  \[
  E(\lambda, S^{-1}TS) = \{ S^{-1}u : u \in E(\lambda, T) \}, \qquad
  E(\lambda, T) = \{ Sv : v \in E(\lambda, S^{-1}TS) \}.
  \]
\end{enumerate}

\end{solutionBox}

\ProblemDivider{5}

\begin{problemBox}
\ProblemAnnounce
\Exercise{5.A.14}
Give an example of an operator on \(\R^4\) that has no (real) eigenvalues.
\end{problemBox}

\begin{solutionBox}
\SolutionAnnounce
\doublespacing

Define \(T \in \Lmap(\R^4)\) by
\[
T(x_1, x_2, x_3, x_4) = (-x_2, x_1, -x_4, x_3).
\]
This is essentially a double rotation (two independent plane rotations).  

Suppose \(\lambda \in \R\) were an eigenvalue with eigenvector \((x_1,x_2,x_3,x_4) \neq 0\). Then
\[
(-x_2, x_1, -x_4, x_3) = (\lambda x_1, \lambda x_2, \lambda x_3, \lambda x_4).
\]
Equating components gives
\[
\begin{cases}
-x_2 = \lambda x_1,\\
x_1 = \lambda x_2,\\
-x_4 = \lambda x_3,\\
x_3 = \lambda x_4.
\end{cases}
\]
If \(\lambda = 0\), then \(x_1 = x_2 = x_3 = x_4 = 0\), a contradiction. So \(\lambda \neq 0\).  
From the first two equations we obtain \(x_1 = \lambda x_2 = \lambda(-\lambda x_1) = -\lambda^2 x_1\), hence \((1+\lambda^2)x_1 = 0\) and thus \(x_1 = 0\). Then \(x_2 = 0\) as well.  
Similarly, the last two equations give \(x_3 = 0\) and \(x_4 = 0\). Hence no nonzero eigenvector exists. Therefore \(T\) has no real eigenvalues.

\end{solutionBox}
\newpage

\ProblemDivider{6}

\begin{problemBox}
\ProblemAnnounce
Show that if \(L(\R)\) stands for the linear maps \(T:\R \to \R\), then \(L(\R)\) is isomorphic to \(\R\).
\end{problemBox}

\begin{solutionBox}
\SolutionAnnounce
\doublespacing

Define \(F : L(\R) \to \R\) by \(F(T) = T(1)\) for each \(T \in L(\R)\).

\begin{itemize}
  \item \textbf{Linearity:}  
  Let \(T_1, T_2 \in L(\R)\) and \(\alpha \in \R\). Then
  \[
  F(T_1 + T_2) = (T_1 + T_2)(1) = T_1(1) + T_2(1) = F(T_1) + F(T_2),
  \]
  \[
  F(\alpha T_1) = (\alpha T_1)(1) = \alpha T_1(1) = \alpha F(T_1).
  \]
  Hence \(F\) is linear.

  \item \textbf{Injectivity:}  
  Suppose \(F(T_1) = F(T_2)\). Then \(T_1(1) = T_2(1)\). For any \(x \in \R\),
  \[
  T_1(x) = T_1(x \cdot 1) = x T_1(1) = x T_2(1) = T_2(x \cdot 1) = T_2(x).
  \]
  Thus \(T_1 = T_2\), so \(F\) is injective.

  \item \textbf{Surjectivity:}  
  Given \(r \in \R\), define \(T_r : \R \to \R\) by \(T_r(x) = r x\). Clearly \(T_r \in L(\R)\) and
  \[
  F(T_r) = T_r(1) = r.
  \]
  Hence \(F\) is surjective.
\end{itemize}

Therefore \(F\) is a vector‑space isomorphism, and \(L(\R) \cong \R\).

\end{solutionBox}

\end{document}
